\section{A uniform sparse-pin estimate}\label{sec:bootstrap}

We now prove that a prescribed sparse set of slot points is unlikely to
belong to cycles avoiding the pins.  This supplies the regular slices
needed for diffusion.  The argument is a finite downward induction on
the pin set: the estimate for a larger pin set controls the geometry
used to remove some of those pins.  It adapts the induction and
obstacle estimates of \cite[Sections~5 and~7]{Ordering}, with the
interaction rates and the coset coordinates retained throughout.

All constants in this section may depend on the fixed parameters of
Section~\ref{sec:geometry}, but not on the cube, pin sets, time grid,
cut, spatial anchor, or inverse temperature.  Dependence on an integer
\(K\) will always be indicated.  A bound \(o_K(a(\beta))\) means that,
for each fixed \(K\), the supremum of the absolute error divided by
\(a(\beta)\) over all these other data tends to zero as
\(\beta\to\infty\).  For an expectation or probability the supremum
is taken after averaging.  We require no uniform rate of convergence
as \(K\) varies.

\subsection{Blocking probabilities from an induction hypothesis}

Set \(p=\beta^{-9/10}\), initially with \(\beta\ge1\), and choose
a finite circular grid \(\mathcal G\subset\mathbb R/\beta\mathbb Z\)
containing zero whose successive gaps are at most \(p^{12}\).
Write \(\mathcal V_{\mathcal G}=V_\Lambda\times\mathcal G\).
For \(P=D_0\cup F\), where \(F\subseteq\mathcal V_{\mathcal G}\),
let \(\mathsf I(P)\) denote the assertion
\begin{equation}\label{eq:induction-assertion}
 \mathbb P_P(T'\subseteq\mathcal B_P)\le p^{|T'|}
 \quad\text{for every nonempty sparse }T'
 \text{ at one grid time, disjoint from }P.
\end{equation}
Here and below sparseness counts slot multiplicities as in
Definition~\ref{def:sparse}.  Until the final subsection we assume
\(\mathsf I(P)\); no estimate for a smaller pin set is being assumed.

Expose \(\mathcal B_P\) under \(\mathbb P_P\).  In a fixed coset
\(\gamma\), let \(\mathcal O_t^\gamma\) be the blocked set in
bar coordinates, as in Definition~\ref{def:regular}.  At a grid
time \(q\), any prescribed sparse set \(A\) of distinct sites in
that coordinate lattice satisfies
\begin{equation}\label{eq:grid-blocking}
 \mathbb P_P(A\subseteq\mathcal O_q^\gamma)\le p^{|A|}.
\end{equation}
To see this, select a fixed slot at each corresponding physical site.
If any selected slot is pinned, or any site lies outside the cube,
the blocking event is impossible.  Otherwise the selected slot
points form a permissible test in \eqref{eq:induction-assertion}.

\begin{lemma}[Blocking at deterministic times]\label{lem:blocking}
Assume \(\mathsf I(P)\).  At every deterministic time \(t\), in
every coset, a prescribed set \(A\) of distinct coordinate sites
satisfies
\begin{align}
 \mathbb P_P(A\subseteq\mathcal O_t^\gamma)&\le Cp^6
 &&\text{if }|A|=6,\label{eq:six-blocking}\\
 \mathbb P_P(A\subseteq\mathcal O_t^\gamma)&\le(C\sqrt p)^{|A|}
 &&\text{if }A\ne\varnothing\text{ is sparse}.
 \label{eq:sparse-blocking}
\end{align}
The constants are uniform in the choice and size of the grid.
\end{lemma}

\begin{proof}
Let \(q\) be the grid time preceding \(t\), taking \(q=t\) when
\(t\in\mathcal G\).  The circular interval \((q,t]\) has length
at most \(p^{12}\).  A site blocked at \(t\) was either blocked
at \(q\) or has an incident inter-site mark in this interval.
The seam does not change the number of holes at any site.  The
total base intensity of all slot edges incident to a site is bounded
by a constant depending only on \(S,J\).  Thus Lemma~\ref{lem:factorial}
bounds the probability of any mark incident to six prescribed sites
by \(Cp^{12}\).  Since a six-site set is sparse, \(L\ge6\) and
\eqref{eq:grid-blocking} give \(p^6+Cp^{12}\le Cp^6\).

For a sparse set of \(n\ge1\) sites, either at least
\(\lceil n/2\rceil\) were blocked at \(q\), or at least that
many have incident marks.  The first event costs at most
\(2^np^{\lceil n/2\rceil}\), by \eqref{eq:grid-blocking} for the
subsets.  In the second event there are at least
\(j=\lceil n/4\rceil\) distinct marks in the union of the edge
stars: one mark is incident to at most two of the selected sites.
If \(N\) counts the marks in that union, factorial domination gives
\[
 \mathbb P_P(N\ge j)
 \le\mathbb E_P\binom Nj
 \le\frac{(Cnp^{12})^j}{j!}
 \le(C'p^{12})^j.
\]
Here \(n/j\le4\) and \(j!\ge(j/e)^j\).  Both event bounds are
at most \((C\sqrt p)^n\) after enlarging \(C\).  This argument
does not require independence between blocking and intervening marks,
nor any bipartite property of the interaction graph.
\end{proof}

\subsection{Sparse witnesses for large blocked components}

We next turn the joint blocking bound into a tail for a connected
obstacle.  Components \(\mathcal C_t^\gamma(w)\) use steps of
supremum length at most eight in the coset coordinates.  Write
\(\operatorname{rad}_w\mathcal C_t^\gamma(w)
=\max_{v\in\mathcal C_t^\gamma(w)}|v-w|_\infty\).
The witness construction below is the one underlying
\cite[Lemma~5.2]{Ordering}; we include the sparseness and counting
arguments because only sparse joint tests are available.

\begin{lemma}[Component tail]\label{lem:component-tail}
There are \(\beta_0,c,C,\lambda>0\) such that, if
\(\beta\ge\beta_0\) and \(\mathsf I(P)\) holds, then
\begin{equation}\label{eq:component-tail}
 \mathbb P_P\bigl(w\in\mathcal O_t^\gamma,
   \operatorname{rad}_w\mathcal C_t^\gamma(w)\ge h\bigr)
 \le C\exp(-ch^\lambda),\qquad h\ge1.
\end{equation}
This holds for every deterministic time \(t\), coset \(\gamma\),
and coordinate site \(w\), with the same constants for all
\(\beta\ge\beta_0\).
\end{lemma}

\begin{proof}
Fix an integer \(b>100\) so large that \(b^\alpha>4\), and set
\(h_j=16b^j\) for \(j\ge0\).  An admissible depth-\(j\) witness
rooted at \(w\) assigns a coordinate site \(x_v\) to each vertex
of the complete rooted ordered binary tree with every leaf at depth
\(j\), with
\(x_\varnothing=w\).  At a vertex having \(k\ge1\) levels
below it, require that the first child has the same site as its
parent, that the second child's site is within distance \(h_k+8\)
of the parent, and that the leaf sets under the two children have
mutual distance at least \(h_k/4\).  Only the leaves are required
to be blocked.  In depth zero the sole leaf is the root.

A blocked path from \(w\) to distance at least \(h_j\) supplies
such a witness.  For \(j\ge1\), truncate the path at its first
site at distance at least \(h_j\).  The first child crossing
starts at \(w\); the second starts at the first path site at
distance at least \(h_j/2\) from \(w\).  That second start is
at distance at most \(h_j/2+8\), and the remaining path reaches
distance at least \(h_j/2-8\) from it.  Each crossing can therefore
be truncated on first reaching distance \(h_{j-1}\) from its
own start.  These truncated crossings lie in the two start-centered
balls of radius \(h_{j-1}+8\), whose mutual distance is at least
\[
 h_j/2-2(h_{j-1}+8)\ge h_j/4.
\]
Recursing inside these crossings constructs the children and keeps
all their descendant leaves in the indicated balls.

Every admissible witness has \(2^j\) distinct leaves and its leaf
set is sparse.  Indeed, a ball of radius \(h<h_k/8\) cannot meet
both branches of any split of scale at least \(h_k\); hence it
contains at most \(2^{k-1}\) leaves.  For
\(h_{k-1}/8\le h<h_k/8\), with \(1\le k\le j\), this is
compatible with the sparse bound because
\[
 L(1+h)^\alpha\ge L(2b^{k-1})^\alpha\ge2^{k-1}.
\]
Below \(h_0/8\) there is at most one leaf.  At and above
\(h_j/8\), the same calculation allows all \(2^j\) leaves.
The depth-zero case is immediate.  These observations cover every
center and every real radius, including the scale endpoints.

Let \(A_j\) bound the number of admissible assignments rooted at
a specified site.  Ignoring restrictions between the two subtrees
gives
\[
 A_0=1,\qquad A_j\le Ch_j^3 A_{j-1}^2,
 \qquad
 2^{-j}\log A_j
 \le\sum_{k=1}^j2^{-k}\log(Ch_k^3)\le\log C_b.
\]
Thus \(A_j\le C_b^{2^j}\).  This numerical overcount does not
declare the unrestricted assignments admissible: only the separated,
and therefore sparse, witnesses enter the probability union bound.
For each such witness, \eqref{eq:sparse-blocking} bounds the
probability of its blocked leaves by \((C\sqrt p)^{2^j}\).
Choose \(\beta_0\) so large that \(C_bC\sqrt p\le e^{-2}\)
for all \(\beta\ge\beta_0\).  Then
\[
 \mathbb P_P\bigl(w\in\mathcal O_t^\gamma,
   \operatorname{rad}_w\mathcal C_t^\gamma(w)\ge h_j\bigr)
 \le e^{-2\cdot2^j}.
\]
For \(h_j\le h<h_{j+1}\), we have
\(2^j\ge(h/(16b))^\lambda\), where
\(\lambda=\log_b2>0\).  Increasing \(C\) for \(1\le h<16\)
proves \eqref{eq:component-tail}.
\end{proof}

\begin{proposition}[Probability of an irregular slice]\label{prop:irregularity}
Under \(\mathsf I(P)\), for \(\beta\ge\beta_0\), every
deterministic time \(t\), anchor \(z\in\mathbb Z^3\), and
radius \(R\ge1\) satisfy
\begin{equation}\label{eq:irregularity}
 \mathbb P_P(t\text{ is not regular for }z\text{ at radius }R)
 \le\delta_R:=CR^3p^6+C\exp(-cR^\xi),
 \qquad \xi=\nu\lambda>0.
\end{equation}
\end{proposition}

\begin{proof}
For each coset there are \(O(R^3)\) centers in \(B_{11R}(z)\).
A ball of the fixed radius \(M_*\) has a fixed number of six-site
subsets.  Equation~\eqref{eq:six-blocking} and a union bound give
\(CR^3p^6\) for failure of the first regularity condition.  Failure
of the second is bounded by Lemma~\ref{lem:component-tail} and
\[
 C\sum_{w:|w-z|_\infty>10R}
  e^{-c(1+|w-z|_\infty)^{\nu\lambda}}
 \le C e^{-c'R^{\nu\lambda}}.
\]
For the last inequality, count at most \(C(1+n)^2\) sites at
distance \(n\) and absorb both this polynomial and the tail sum
into a smaller exponential rate.  Finally sum over the finitely
many cosets.
\end{proof}

We have obtained regularity at prescribed times using only the
induction hypothesis for \(P\).  To use this information over a
walk's whole history, fix an integer \(K\ge1\) and set
\begin{equation}\label{eq:scales}
 \tau=K\beta,\qquad R=R_K=\lceil a(K+1)\beta\rceil.
\end{equation}
The constant \(a\ge1\) will be fixed sufficiently large in
Lemma~\ref{lem:localization}, independently of \(K\).  We also
take \(a\ge R_0\), the fixed radius threshold from
Section~\ref{sec:geometry}.  Then
\begin{equation}\label{eq:irregularity-scale}
 \delta_{R_K}=O_K(\beta^{-12/5})=o_K(\beta^{-3/2}).
\end{equation}
Call the exposure \emph{good for \(z,R\)} if its nonregular times
in one period have total Lebesgue measure at most \(\beta/100\).
Denote this event by \(G_{z,R}\).  Fubini's theorem and Markov's
inequality give
\begin{equation}\label{eq:good-exposure}
 \mathbb P_P(G_{z,R}^c)\le100\delta_R.
\end{equation}
Neither this assertion nor Proposition~\ref{prop:irregularity}
requires independence between time slices.

\subsection{Localization and diffusion}

To localize at \(z\), kill the walk whenever its site leaves
\(\mathfrak B_R(z)\), in addition to the boundary killing already
present in \(U\).  Apply the killing projections also at starting
and ending layers and on the appropriate sides of instantaneous
operations.  Localized kernels, extended by zero to the original
hole layers, are entrywise smaller and retain the norm contractions
of Section~\ref{sec:geometry}.  Write \(U_{d,R}\) for a localized
full traversal in direction \(d\in\{+,-\}\); reversing all factors
gives \(U_{-,R}=U_{+,R}^*\).

\begin{lemma}[Localization]\label{lem:localization}
One may choose \(a\) in \eqref{eq:scales} independently of \(K\)
so that the following holds, without assuming \(\mathsf I(P)\).
Let \(i\) be a deterministic slot point that is almost surely a
hole under the exposure law for an allowed pin set \(P\), and put
\(z=\bar x_i\).  Let \(\mathcal K\) be a deterministic schedule
of at most \(2K+2\) full or partial boundary-killed traversals
starting at \(i\), each in either direction and each covering at
most one period.  Use matching layers at every concatenation.
If \(\mathcal K_R\) is localized throughout at \(z\), then
\begin{equation}\label{eq:localization-loss}
 \mathbb E_P\bigl\|(\mathcal K-\mathcal K_R)
          \ind{\{i\}}\bigr\|_1\le C_K e^{-c\beta}.
\end{equation}
The same estimate holds for localization of the walk on the full
slot lattice, with all slots available and no exposed trajectories.
\end{lemma}

\begin{proof}
Sample the actual picture under \(\mathbb P_P\), thereby giving
the exposure its correct marginal.  Independently supply fresh
Poisson clocks for each traversal piece.  Accept these clocks where
both endpoints are holes.  Forced transports use actual picture
marks; fresh uniform completions give the seam transitions.  This
generates the fresh walk conditional on the exposure, in either
direction.  Its expected mass loss from localization is at most
its probability of exiting \(\mathfrak B_R(z)\).

An exit requires at least \(R/B_*\) inter-site steps.  Before the
first exit, one of the pieces must therefore make at least
\[
 j=\left\lfloor\frac{R}{B_*(2K+2)}\right\rfloor
 \ge\frac{a\beta}{2B_*}-1
\]
steps, starting at a site in \(\mathfrak B_R(z)\).  Within a
single piece the actual marks in this list are distinct: physical
time runs injectively through at most one period, also for a reverse
piece or a piece crossing the seam.  The fresh marks are likewise
distinct, and seam completions cause no inter-site step.

For a fixed starting site, ignore slot compatibility to overcount
all lists of \(j\) steps.  There are a bounded number of possible
next sites, slot edges, and actual-or-fresh designations at each
step.  Lemma~\ref{lem:factorial}, combined with independence of
the fresh Poisson family, bounds the mixed factorial measure by
the product of the actual intensity multiplied by two and the
fresh intensity.  On unwrapping the chosen piece, the chronological
times range over a simplex of volume at most \(\beta^j/j!\).
After summing edge choices and designations, the expected number
of lists is consequently at most
\[
 \frac{(C\beta)^j}{j!}.
\]
Choose \(a\) so large that \(j\ge a\beta/(3B_*)\) for large
\(\beta\) makes this quantity at most \(e^{-c\beta}\), by
\(j!\ge(j/e)^j\).  Sum over the at most \(2K+2\) pieces and
the \(O(R^3)\) possible starting sites; the polynomial factor is
absorbed by reducing \(c\) and using a constant \(C_K\).

All factorial estimates were applied unconditionally under
\(\mathbb P_P\), within a single piece at a time.  Actual marks
may be reused on later traversals; no independence of those uses
was asserted.  For the full-lattice walk the same argument uses
only fresh marks, and gives the final assertion.
\end{proof}

\begin{proposition}[Diffusion on a good exposure]\label{prop:diffusion}
Fix an exposure good for \(z,R_K\).  Let a kernel be a subinterval
of a concatenation of full traversals in arbitrary directions,
localized throughout at \(z\).  For every nonnegative input of
mass at most one, its output after elapsed time \(b'\ge\beta/16\)
satisfies
\begin{equation}\label{eq:diffusion-l2}
 \|v_{\rm out}\|_2\le C(b')^{-3/4}.
\end{equation}
The same statement holds for the adjoint of any such subinterval.
For such a kernel \(\mathcal K_R\) of length \(b'\ge\beta/8\),
\begin{equation}\label{eq:diffusion-entry}
 \max_{j,k}(\mathcal K_R)_{jk}\le C(b')^{-3/2}.
\end{equation}
The constants are independent of \(K\) and \(R_K\).
\end{proposition}

\begin{proof}
This is Nash's dissipation argument, applied only during regular
times; see \cite[pp.~935--936]{Nash1958}.  Let
\(q=\|v\|_2^2\) for an evolving localized mass.  At a regular
smooth time, Proposition~\ref{prop:nash} and
\eqref{eq:norm-dissipation} imply
\[
 q'\le-cq^{5/3},\qquad (q^{-2/3})'\ge c'>0
 \quad\text{when }q>0.
\]
At other smooth times and at instantaneous operations the norm
does not increase.  If it becomes zero, the conclusion is immediate.
The input and output masses are at most one, by contraction.

The subinterval has only two possible partial end pieces; all
intermediate pieces have length \(\beta\).  Its total nonregular
time is therefore at most \((b'/\beta+2)\beta/100\).  For
\(b'\ge\beta/16\), its regular time is at least \(67b'/100\).
Integrating the displayed differential inequality proves
\eqref{eq:diffusion-l2}.  The same count holds after reversing the
subinterval, and the reversed instantaneous factors are contractions.
For \eqref{eq:diffusion-entry}, split at the temporal midpoint:
an entry is the inner product of a forward mass from its starting
slot and an adjoint mass from its ending slot.  Each propagates
for \(b'/2\ge\beta/16\), so Cauchy--Schwarz and
\eqref{eq:diffusion-l2} give the result.
\end{proof}

\subsection{The hitting energy and the induction step}

We have proved diffusion under the geometry supplied by
\(\mathsf I(P)\).  It remains to turn this into the assertion
for \(E\) when \(P=E\cup T\).  The following two energy identities
separate the return event from the spatial estimates.  They hold at
each fixed exposure and do not require the induction hypothesis.

\begin{lemma}[Hitting energy]\label{lem:hitting-energy}
Let \(E\) be allowed and let \(T\ne\varnothing\) be a finite
set of slot points at a deterministic cut \(\theta\), disjoint
from \(E\).  Put \(P=E\cup T\), \(m_T=|T|\), and
\(\eta=\ind{T}\).  At almost every exposure under \(\mathbb P_P\),
define \(f=1\) on \(T\); at every other cut hole let \(f\) be
the probability that the boundary-killed fresh walk ever hits
\(T\) at a subsequent integer turn.  Then \(0\le f\le1\),
\(f\) has zero boundary values, and, for \(g=f-\eta\),
\begin{equation}\label{eq:hitting-energy}
 \sum_{i\in T}q_i(\mathbf a)
 \le m_T-\langle f,A_Qf\rangle
 =\langle\eta,Q\eta\rangle+2\langle Q\eta,g\rangle
       -\langle g,A_Qg\rangle.
\end{equation}
\end{lemma}

\begin{proof}
As in \cite[Lemma~4.5]{Ordering}, construct \(f\) from finite
hitting horizons.  Set \(f_0=\eta\), and define
\(f_{n+1}=1\) on \(T\) and \(f_{n+1}=U^*f_n\) off \(T\).
Positivity and substochasticity give
\(0\le f_n\le f_{n+1}\le1\).  Their coordinatewise limit
has the stated hitting interpretation and satisfies
\(U^*f=f\) off \(T\), with zero boundary coordinates.
No eventual absorption assumption is needed.

Since \(T\) is present only at the cut, a strictly positive
return to \(T\) occurs at a positive integer turn.  Removing the
extra pins of \(E\setminus D_0\) from the stopping rule yields
\(q_i(\mathbf a)\le(U^*f)_i\) for \(i\in T\).  Harmonicity
off \(T\) and \(f=1\) on \(T\) give
\[
 \langle f,A_Qf\rangle
 =\langle(I-U^*)f,f\rangle
 =m_T-\sum_{i\in T}(U^*f)_i.
\]
The quadratic forms of \(U\), \(U^*\), and \(Q\) agree on
real vectors.  Finally, \(\langle\eta,g\rangle=0\); expanding
\(f=\eta+g\) gives \eqref{eq:hitting-energy}.
\end{proof}

For an arbitrary real vector \(y\) on the cut holes with zero
boundary values, set \(E_y=\langle y,A_Qy\rangle\).  For either
direction \(d\in\{+,-\}\), split the traversal at elapsed time
\(u\in I=[\beta/3,2\beta/3]\) as
\[
 U_d=B_d(u)C_d(u),\qquad Y_d(u)=B_d(u)^*y.
\]
Here \(C_d(u)\) is the starting portion and \(B_d(u)\) the
ending portion; both factors use matching hole layers at physical
time \(t_d(u)=\theta+u\) or \(\theta-u\) modulo \(\beta\).
Each instantaneous operation is assigned to precisely one factor.
The exceptional split times do not affect the integrals below.

\begin{lemma}[Suffix energy]\label{lem:suffix-energy}
With these definitions and zero extension to the available-slot
graph at the split,
\begin{equation}\label{eq:suffix-energy}
 \int_I\mathcal D_{t_d(u)}(Y_d(u))\,du
 \le\frac12\bigl(\|y\|_2^2-\|U_d^*y\|_2^2\bigr)
 \le E_y.
\end{equation}
Moreover \(\|Y_d(u)\|_\infty\le\|y\|_\infty\).
\end{lemma}

\begin{proof}
As \(u\) decreases, \(B_d(u)^*y\) applies successively more
adjoint factors from the endpoint.  Equation~\eqref{eq:norm-dissipation}
identifies its continuous squared norm loss with twice the displayed
Dirichlet energy; instantaneous factors cannot increase the norm.
Integrating over the subinterval \(I\) gives the first inequality.
The second follows from the exact identity
\[
 E_y-\tfrac12(\|y\|_2^2-\|U_d^*y\|_2^2)
 =\tfrac12\|y-U_d^*y\|_2^2\ge0,
\]
since \(\langle y,U_d^*y\rangle=\langle y,Qy\rangle\).
Supremum norm contraction proves the final assertion.
\end{proof}

\begin{theorem}[Uniform sparse-pin estimate]\label{thm:sparse-pins}
There is \(\beta_1<\infty\), independent of the cube and pins,
such that for every \(\beta\ge\beta_1\), every allowed pin set
\(E\), and every nonempty sparse set \(T\) at one deterministic
time, disjoint from \(E\),
\begin{equation}\label{eq:sparse-pins}
 \mathbb P_E(T\subseteq\mathcal B_E)
 \le\beta^{-9|T|/10}.
\end{equation}
Consequently \eqref{eq:irregularity} and \eqref{eq:good-exposure}
hold under every allowed pin law, with no induction hypothesis.
\end{theorem}

\begin{proof}
First fix a grid \(\mathcal G\) as above and prove
\(\mathsf I(E)\) by downward induction on the finite number
of grid pins outside \(D_0\).  The assertion is vacuous when
all grid points are pinned.  Suppose it is proved for every
strict superset of \(E\) among these pin sets, and fix an
admissible sparse test \(T\) for \(E\).  Then
\(P=E\cup T\) satisfies \(\mathsf I(P)\), so all the preceding
probability estimates apply under \(\mathbb P_P\).
Use \(K=1\), \(R=R_1\), and the vectors of
Lemma~\ref{lem:hitting-energy}.

The first term in \eqref{eq:hitting-energy} obeys
\begin{equation}\label{eq:bootstrap-short-return}
 \mathbb E_P\langle\eta,Q\eta\rangle
 \le Cm_T\bigl(e^{-c\beta}+\delta_R
                       +R^\alpha\beta^{-3/2}\bigr).
\end{equation}
Indeed, propagate a unit mass from each \(i\in T\) in either
direction, localizing at \(z_i=\bar x_i\).  Lemma~\ref{lem:localization}
bounds the expected discarded mass by \(Ce^{-c\beta}\).
On an exposure good for \(z_i,R\), Proposition~\ref{prop:diffusion}
bounds each one-period entry by \(C\beta^{-3/2}\), and
sparseness allows at most \(L(1+R)^\alpha\) target points of
\(T\) in that region.  On a bad exposure the total output
mass is at most one, and its probability is at most \(100\delta_R\).
Sum over \(i\) and average the directions.

It remains to absorb the linear term by its accompanying energy.
Write \(E_g=\langle g,A_Qg\rangle\), use the splits above with
\(y=g\), and let \(C_{d,R,z_i}(u)\) denote the starting
portion localized at \(z_i\).  Define a mass on the split layer by
\[
 M_d(u)=\sum_{i\in T}
  \ind{\{t_d(u)\text{ regular for }z_i\text{ at radius }R\}}
  C_{d,R,z_i}(u)\ind{\{i\}},
 \qquad
 F_d=\frac1{|I|}\int_I\langle M_d(u),Y_d(u)\rangle\,du.
\]
The unchanged pairing
\(\langle U_d\eta,g\rangle
=\langle C_d(u)\eta,Y_d(u)\rangle\) is constant in \(u\).
Since \(0\le g\le1\), the suffix has supremum norm at most one.
Localization, deterministic-time regularity, and Fubini therefore give
\begin{equation}\label{eq:bootstrap-pairing-error}
 \mathbb E_P\bigl|\langle U_d\eta,g\rangle-F_d\bigr|
 \le Cm_T(e^{-c\beta}+\delta_R).
\end{equation}
No good-exposure indicator is needed here: the next estimate uses
only regularity at the split slice.

Each retained summand of \(M_d(u)\) has mass at most one and
is supported above \(\mathfrak B_R(z_i)\); the list of centers,
including repetitions, has the growth bound of \(T\).
Proposition~\ref{prop:sparse-energy}, followed by Cauchy--Schwarz
in \(u\) and Lemma~\ref{lem:suffix-energy}, gives pointwise
\[
 |F_d|\le
 \left(\frac{Cm_TR^\alpha}{|I|}E_g\right)^{1/2}.
\]
Because \(|I|=\beta/3\) and \(E_g\ge0\), completing the square
in \(\sqrt{E_g}\) yields
\[
 F_++F_--E_g\le\frac{C m_T R^\alpha}{\beta}.
\]
Now \(2\langle Q\eta,g\rangle
=\langle U_+\eta,g\rangle+\langle U_-\eta,g\rangle\).
Use the fresh-return identity \eqref{eq:fresh-return},
\eqref{eq:hitting-energy}, and
\eqref{eq:bootstrap-short-return}--\eqref{eq:bootstrap-pairing-error}
to conclude
\begin{equation}\label{eq:bootstrap-return}
 r(E,T)\le C\bigl(e^{-c\beta}+\delta_R
              +R^\alpha\beta^{-3/2}+R^\alpha\beta^{-1}\bigr)
 =o(\beta^{-9/10}).
\end{equation}
Here \(R=\lceil2a\beta\rceil\), so the largest polynomial
term is \(O(\beta^{-99/100})\); the error is uniform in every
induction datum.  Increase \(\beta_1\), independently of the
grid and its size, so that \(2r(E,T)\le p\).
Proposition~\ref{prop:pin-identities} now gives
\(\mathbb P_E(T\subseteq\mathcal B_E)\le(2r(E,T))^{m_T}
\le p^{m_T}\), completing the induction.

For arbitrary allowed \(E\) and a test at any deterministic time,
choose a finite grid containing zero, that test time, and all the
finitely many extra-pin times, then refine it to mesh at most
\(p^{12}\).  The uniform grid result proves
\eqref{eq:sparse-pins}.  The same choice of grid for an arbitrary
allowed \(P\) removes the hypothesis \(\mathsf I(P)\) from the
regularity and good-exposure bounds.
\end{proof}

All time integrals above are of measurable functions of the exposure.
The histories have finitely many jumps, their transition kernels are
measurable, and the hitting function is an increasing limit of
finite-horizon probabilities.  The sharp one-slot estimate of order \(\beta^{-3/2}\) has not been
used in this induction; it will follow from the refined energy
argument in Section~\ref{sec:returns}.
